\section{Results}
\label{sec:results}

\subsection{The price of H2}
\label{sec:price}

\begin{table}[t]
\centering
\small
\begin{tabular}{@{}lrrrc@{}}
\toprule
\textbf{scheme} & \textbf{\Sz{}} & \textbf{\So{}} & $\So/\Sz$ & \textbf{slots occupied} \\
\midrule
Strong-4D-Var & 0.8116 & 1.4617 & \textbf{1.801} & $\Psimean$ \\
ETKF          & 0.8883 & 1.4998 & \textbf{1.688} & $\Psimean + \PsiG$ \\
EnKF          & 0.9131 & 1.5381 & \textbf{1.684} & $\Psimean + \PsiG$ \\
\addlinespace
SDA3          & 0.5365 & 0.5355 & 0.998 & all three \\
DirectUNet-M  & 0.5064 & 0.5072 & 1.002 & $\Psimean$ \\
CFM-M         & 0.4810 & 0.4784 & 0.995 & all three \\
DirectUNet\,$+$\,SDA3 & \textbf{0.4204} & 0.4183 & 0.995 & composed \\
\bottomrule
\end{tabular}
\caption{Lorenz-96, 200 windows, pooled RMSE. \textbf{The two blocks are
different method classes} and the lower ratios must be read through
Table~\ref{tab:classes}: DirectUNet and CFM are model-free, so their
$\So/\Sz \approx 1$ is definitional. Only the classical block's degradation
measures the cost of model error. Source:
\texttt{reports/l96/outputs/l96\_consolidated\_benchmark.md}.}
\label{tab:price}
\end{table}

Classical DA degrades by $1.68$--$1.80\times$ from \Sz{} to \So{}. That is the
unbuffered cost of H2: by~\eqref{eq:ssm} the model \emph{is} the prior, so a
parameter bias is not a perturbed input but a corrupted prior, compounding over
a 500-step window.

\subsection{Weak-constraint 4D-Var, the formulation that relaxes H2}
\label{sec:weak4dvar}

Lorenz-63 is small enough to afford the weak-constraint formulation, which frees
a per-step model-error term and is therefore the classical method permitted to
know that the model is wrong.

\begin{table}[t]
\centering
\small
\begin{tabular}{@{}lrr@{}}
\toprule
\textbf{scheme} & \textbf{\Sz{}} & \textbf{\So{}} \\
\midrule
Weak-4D-Var   & \textbf{0.642} & \textbf{1.633} \\
Strong-4D-Var & 0.767 & 2.100 \\
EnKF          & 0.881 & 2.267 \\
ETKF          & 0.878 & 2.271 \\
\bottomrule
\end{tabular}
\caption{Lorenz-63 state RMSE. Weak-4D-Var is best in both scenarios, and under
model error it beats strong-constraint by 22\%.
Source: \texttt{reports/l63/outputs/s0\_s1\_synthesis.md}.}
\label{tab:weak4dvar}
\end{table}

This is the honest ceiling for the classical family, and it sharpens rather than
weakens the argument. Relaxing H2 \emph{within} the classical framework buys a
great deal --- so the residual gap to the learned schemes is not an artefact of
comparing against a strawman. It also bounds what C1 may claim: the
marginal-value collapse reported in Section~\ref{sec:marginal} is established for
strong-constraint and filtering methods, and the corresponding weak-constraint
measurement on Lorenz-96 does not yet exist.
\needsrun{Weak-constraint 4D-Var on Lorenz-96 across the observation-density
grid. The implementation exists; the benchmark run does not.}

\subsection{Conditioning on model information is nearly inert}
\label{sec:inert}

If the learned schemes are robust to \So{}, the obvious question is whether they
use $\thv^{*}$ at all.

\begin{table}[t]
\centering
\small
\begin{tabular}{@{}llrrr@{}}
\toprule
\textbf{scheme} & \textbf{conditioning} & \textbf{\Sz{}} & \textbf{\So{}} & $\So/\Sz$ \\
\midrule
SDA1 & none (unconditional prior) & 0.5532 & 0.5521 & 0.998 \\
SDA2 & \emph{true} $\thv$ $+$ corrupted $\zv^{*}$ & 0.5588 & 0.5610 & 1.004 \\
SDA3 & \emph{noisy} $\thv^{*}$ & \textbf{0.5365} & 0.5355 & 0.998 \\
\bottomrule
\end{tabular}
\caption{Lorenz-96 conditioning ladder, pooled RMSE. Conditioning on \emph{true}
parameters is 1.0\% \emph{worse} than not conditioning at all; conditioning on
\emph{noisy} ones is 3.0\% better.}
\label{tab:ladder}
\end{table}

Across three independent families the effect of model-information conditioning
on \Sz{} skill is at most $\sim$1.5\%, and \emph{the sign is not consistent}: the
SDA family gives $+1.0\%$ (worse) for true parameters and $-3.0\%$ (better) for
noisy ones, a non-monai SDA variant gives $-1.5\%$, and a $\tau{=}0$ CFM
conditioned on corrupted forcing gives $+0.9\%$ (worse). Conditioning is, to a
good approximation, \emph{inert}; the one variant that reliably helps is the
noisy one, which looks like regularisation rather than information use.

This is the sharpest available statement of the thesis. The learned schemes are
robust to corrupted model parameters largely \emph{because they barely depend on
them} --- the observations carry the information. Classical DA has no such
fallback: its prior \emph{is} the integrated model. H2 binds not because model
information is unhelpful, but because DA cannot down-weight it.

\paragraph{The training-exposure control has been run.}
The configuration trained with zero forcing/state bias --- the only one never
exposed to \So{}-level error in training --- gives \Sz{} 0.7046 / \So{} 0.7033,
ratio 0.998, indistinguishable from its exposed sibling (0.997). Removing
training-time exposure entirely leaves the invariance intact.
\caveat{This control has no counterpart in the canonical benchmark table, which
is scoped to one backbone family; it is cited to its own run artefact.}

\paragraph{The same ladder on QG.} The QG DirectUNet variants repeat the pattern
at a different scale and resolution: the oracle-conditioned variant attains the
best \Sz{} explained variance (0.9827) but falls to 0.9490 under \So{}, while the
noisy-conditioned variants hold (0.9798 $\to$ 0.9763, and 0.9824 $\to$ 0.9808).
Conditioning on \emph{true} model information is again the fragile choice.

\subsection{ODE sensitivity is not posterior sensitivity}
\label{sec:sensitivity}

Two distinct sensitivities are in play, and classical DA cannot separate them:
the \emph{model} sensitivity $\partial \xv_{\mathrm{ODE}}/\partial(\thv,\zv)$ of
the free-running trajectory, Lyapunov-amplified and large; and the
\emph{posterior} sensitivity
$\partial\, p(\xv_1\mid\Cv)/\partial(\thv,\zv)$ at fixed $\yv$.

Classical DA conflates them \textbf{by construction}, because its estimator
\emph{is} the model: by~\eqref{eq:ssm} every route from $(\thv,\zv)$ to the
analysis passes through $\Phi$. Its \Sz{}$\to$\So{} degradation therefore measures
\emph{its own estimator's} sensitivity, and cannot be read as a measurement of
how much $(\thv,\zv)$ uncertainty matters for the posterior.

A conditioned amortised estimator measures the other quantity directly: it takes
$(\thv^{*},\zv^{*})$ as inputs at fixed $\yv$, so perturbing the conditioning
while holding observations fixed is a finite-difference estimate of posterior
sensitivity. That perturbation moves skill by $\leq 1.5\%$
(Section~\ref{sec:inert}), and full corruption by $\sim 0.2\%$.

\paragraph{Consequence.}
Assessing parameter and forcing uncertainty by \emph{propagating it through the
model} systematically \textbf{overstates} its effect on the posterior whenever
the observations are informative, and the overstatement is exactly the gap
between the two sensitivities --- here $1.68$--$1.80\times$ against $\leq 1.5\%$.

\subsection{Model error destroys the marginal value of observations}
\label{sec:marginal}

\begin{table}[t]
\centering
\small
\begin{tabular}{@{}lccc@{}}
\toprule
\textbf{method} & \textbf{\Sz{}: gain from $2\times$ obs} & \textbf{\So{}: same gain} & \textbf{collapse} \\
\midrule
Strong-4D-Var & $0.9701 \to 0.7788$ (\textbf{19.7\%}) & $1.4751 \to 1.4276$ (\textbf{3.2\%}) & $\mathbf{6.2\times}$ \\
ETKF          & $1.0973 \to 0.8815$ (19.7\%)          & $1.6367 \to 1.4680$ (10.3\%)         & $1.9\times$ \\
EnKF          & $1.0927 \to 0.9046$ (17.2\%)          & $1.6503 \to 1.5022$ (9.0\%)          & $1.9\times$ \\
\bottomrule
\end{tabular}
\caption{Lorenz-96, \emph{temporal} observation density doubled (15 $\to$ 30
observations per window). A separate configuration from Table~\ref{tab:price}
--- all five parameters randomised $\pm 20\%$ --- so the two must not be
cross-quoted. Source:
\texttt{reports/l96/outputs/s0\_s1\_obs\_density\_da\_baselines.md}.}
\label{tab:marginal}
\end{table}

Under model error, doubling the observations buys strong-constraint 4D-Var
3.2\%. In the fast-variable subspace, explained variance goes
$+0.417 \to -0.011$ at the sparser density and $+0.622 \to +0.047$ at the denser
one.

\textbf{H2 does not merely set an error floor; it decouples the analysis from the
observations.} A method whose prior is a wrong model cannot be rescued by more
data --- exactly the regime of sparse-observation geophysical DA. This is also
H3b's shadow: at low density DA has nothing but the model, so model error and
observation sparsity \emph{compound rather than add}.

\subsection{The state variable changes skill at identical physics}
\label{sec:representation}

The QG \texttt{psi\_state} variant integrates the same dynamics with the
streamfunction as the state variable. The $q$-space physics is bit-identical ---
free forecasts agree to $\sim 10^{-6}$ relative --- and only the representation
differs, and with it the observation operator, which becomes a trivial index
lookup.

It gives the \emph{best per-field streamfunction analysis}
($\psi_1,\psi_2$ EV $\approx 0.98$) and a \emph{degenerate potential-vorticity
field} ($q$ EV $=-3.2$ against the $q$-state reference $+0.34$), because
$q \approx \nabla^2\psi$ amplifies high-wavenumber $\psi$-analysis error by
$K^2$. At matched observation noise the $\psi$-state $q$-field EV is 0.583
against the $q$-state 0.752.

Identical dynamics, identical information, and a large skill difference produced
purely by which variable is called ``the state''. The representation is not a
neutral container.
\caveat{The headline $-3.2$ is partly a metric choice --- the aggregate score
reports the $q$ field. The per-field numbers above are the ones to quote.}

\subsection{The non-Gaussian deficit, and what blending buys}
\label{sec:nongaussian}

Table~\ref{tab:zeroing} predicts that the ensemble-Kalman class cannot represent
$\PsiNG$ at any tuning. The measured consequence is that \emph{no scheme
evaluated is simultaneously accurate and calibrated}. The only configuration
carrying a full non-Gaussian branch is 20--25\% over-dispersive --- spread/skill
relative to target $1.20 / 1.22 / 1.24$ at $N = 6/10/30$, stable in $N$ and hence
a sampler property rather than a sample-size artefact --- while every scheme that
closes the RMSE gap does so by importing a deterministic $\Psimean$ and collapses
to ratios $0.28$--$0.59$ against a target of $0.967$.

Blending is where this becomes a design choice rather than an accident.
DirectUNet\,$+$\,SDA3 is the most accurate scheme in Table~\ref{tab:price}
(0.4204, against 0.5064 for its own $\Psimean$ alone and 0.5365 for the sampler
alone), so composing the slots demonstrably buys accuracy. But it buys it by
importing a deterministic mean, which is precisely the move that collapses
dispersion. \textbf{Accuracy and calibration are being traded against each other
through the same mechanism}, and the partition~\eqref{eq:partition} is what makes
that visible: the trade is not between schemes but between which slots are filled
generatively and which are imported.
\needsrun{A calibrated comparison across all classes at one uniform ensemble
convention, with rank histograms. Several cells of the benchmark are still
single-pass point estimates.}

\subsection{A motivating ordering that is not yet evidence}
\label{sec:ordering}

At \Sz{} on Lorenz-96 the schemes order by how much of the window they treat at
once: sequential filters (ETKF 0.8883, EnKF 0.9131) are worse than the window
variational method (Strong-4D-Var 0.8116), which is in turn far worse than the
window-amortised learned schemes (0.4204--0.5365). Read against
Table~\ref{tab:zeroing}, the ordering is not monotone in the number of operator
slots occupied --- EnKF fills $\PsiG$ and still trails a $\Psimean$-only
variational scheme --- which is itself informative: \emph{how much of the problem
is solved jointly} matters here more than which posterior components are
representable.

This is suggestive of H1 but \textbf{confounded}: amortisation, learning and
capacity all vary at once, and the learned schemes see a training set the
classical ones do not.
\needsrun{Filter vs.\ window-smoother vs.\ window-amortised at matched model,
matched observations and matched compute. This is what turns the ordering into
evidence for H1.}
